\chapter{Combinational Arithmetic \& Logic Circuits: Half/Full Adders \& Subtractors}
\label{chap:adders}

\section{Overview and Laboratory Learning Objectives}
Arithmetic Logic Units (ALUs) form the computational core of microprocessors, digital signal processors, and graphics pipelines. The most elementary mathematical operation implemented in hardware is binary addition. This chapter explores the systematic gate-level design, truth-table derivation, hardware breadboarding, and diagnostic verification of Half Adders, Full Adders, Half/Full Subtractors, and 4-bit parallel binary adders.

\begin{important}[Unit Learning Objectives]
Upon completion of this chapter and lab experiments, the student will be proficient in:
\begin{enumerate}[leftmargin=1.5em]
    \item Formulating truth tables for binary addition and binary subtraction.
    \item Deriving minimal sum-of-products Boolean expressions for Sum, Carry, Difference, and Borrow.
    \item Designing and wiring a Half Adder using discrete XOR and AND gates.
    \item Designing and wiring a Full Adder by cascading two Half Adders and an OR gate.
    \item Implementing arithmetic adders using universal NAND gates exclusively.
    \item Analyzing the operation and carry propagation delay of the IC 7483 4-bit binary adder.
    \item Troubleshooting combinational glitches and wiring errors on digital trainer modules.
\end{enumerate}
\end{important}

\section{Binary Addition Principles and the Half Adder}
Binary addition follows four primitive single-bit rules:
\[ 0 + 0 = 0, \quad 0 + 1 = 1, \quad 1 + 0 = 1, \quad 1 + 1 = 10_2 \,\,(\text{Sum } 0 \text{ with Carry } 1) \]

\begin{definition}[Half Adder]
A Half Adder is a combinational circuit with two single-bit inputs ($A, B$) and two outputs: Sum ($S$) and Carry ($C$).
\end{definition}

\begin{center}
\begin{tabular}{cc|cc}
\toprule
\textbf{Input $A$} & \textbf{Input $B$} & \textbf{Sum ($S$)} & \textbf{Carry ($C$)} \\
\midrule
0 & 0 & 0 & 0 \\
0 & 1 & 1 & 0 \\
1 & 0 & 1 & 0 \\
1 & 1 & 0 & 1 \\
\bottomrule
\end{tabular}
\end{center}

Extracting Boolean minterms directly from the truth table:
\begin{keyequation}[Half Adder Boolean Expressions]
\[ S = \overline{A}B + A\overline{B} = A \oplus B \]
\[ C = A \cdot B \]
\end{keyequation}
Hardware implementation requires:
\begin{itemize}
    \item One 2-input XOR gate (1/4th of IC 7486) for Sum.
    \item One 2-input AND gate (1/4th of IC 7408) for Carry.
\end{itemize}

\section{The Full Adder Architecture}
When multi-bit binary numbers are added, each column except the least significant bit must accommodate an incoming carry ($C_{\mathrm{in}}$) generated by the preceding stage.

\begin{definition}[Full Adder]
A Full Adder is a combinational circuit with three single-bit inputs: $A, B$, and $C_{\mathrm{in}}$, and two outputs: Sum ($S$) and Carry Out ($C_{\mathrm{out}}$).
\end{definition}

\begin{center}
\begin{tabular}{ccc|cc}
\toprule
\textbf{Input $A$} & \textbf{Input $B$} & \textbf{Input $C_{\mathrm{in}}$} & \textbf{Sum ($S$)} & \textbf{Carry Out ($C_{\mathrm{out}}$)} \\
\midrule
0 & 0 & 0 & 0 & 0 \\
0 & 0 & 1 & 1 & 0 \\
0 & 1 & 0 & 1 & 0 \\
0 & 1 & 1 & 0 & 1 \\
1 & 0 & 0 & 1 & 0 \\
1 & 0 & 1 & 0 & 1 \\
1 & 1 & 0 & 0 & 1 \\
1 & 1 & 1 & 1 & 1 \\
\bottomrule
\end{tabular}
\end{center}

\subsection{Boolean Derivation for Full Adder}
The Sum output minterms are:
\[ S = \sum m(1, 2, 4, 7) = \overline{A}\overline{B}C_{\mathrm{in}} + \overline{A}B\overline{C_{\mathrm{in}}} + A\overline{B}\overline{C_{\mathrm{in}}} + ABC_{\mathrm{in}} \]
Factoring out $C_{\mathrm{in}}$:
\[ S = C_{\mathrm{in}}(\overline{A}\overline{B} + AB) + \overline{C_{\mathrm{in}}}(\overline{A}B + A\overline{B}) = C_{\mathrm{in}}\overline{(A \oplus B)} + \overline{C_{\mathrm{in}}}(A \oplus B) = A \oplus B \oplus C_{\mathrm{in}} \]
The Carry Out minterms are:
\[ C_{\mathrm{out}} = \sum m(3, 5, 6, 7) = AB + BC_{\mathrm{in}} + AC_{\mathrm{in}} = AB + C_{\mathrm{in}}(A + B) = AB + C_{\mathrm{in}}(A \oplus B) \]

\begin{keyequation}[Full Adder Canonical Expressions]
\[ S = A \oplus B \oplus C_{\mathrm{in}} \]
\[ C_{\mathrm{out}} = AB + C_{\mathrm{in}}(A \oplus B) \]
\end{keyequation}

\subsection{Implementation Using Two Half Adders and an OR Gate}
\begin{enumerate}
    \item \textbf{Half Adder 1}: Inputs $A, B \implies S_1 = A \oplus B$, $C_1 = AB$.
    \item \textbf{Half Adder 2}: Inputs $S_1, C_{\mathrm{in}} \implies S = S_1 \oplus C_{\mathrm{in}} = (A \oplus B) \oplus C_{\mathrm{in}}$, $C_2 = S_1 \cdot C_{\mathrm{in}} = (A \oplus B)C_{\mathrm{in}}$.
    \item \textbf{OR Gate}: $C_{\mathrm{out}} = C_1 + C_2 = AB + (A \oplus B)C_{\mathrm{in}}$.
\end{enumerate}

\section{Laboratory Experiment 5: Breadboarding and Testing Adders}
\begin{experimentbox}[Laboratory Experiment 5 Protocol]
\textbf{Objective}: To design, assemble on a solderless breadboard, and functionally verify the complete truth tables of a Half Adder and a Full Adder.\\
\textbf{Apparatus Required}:
\begin{enumerate}
    \item Digital Trainer Module with $+5\,\mathrm{V}$ supply and logic switches --- 1 unit
    \item IC 7486 (Quad 2-Input XOR) --- 1 piece
    \item IC 7408 (Quad 2-Input AND) --- 1 piece
    \item IC 7432 (Quad 2-Input OR) --- 1 piece
    \item Breadboard and patch jumpers
\end{enumerate}
\end{experimentbox}

\begin{observationtable}[Observation Table 5.1: Full Adder Truth Table Verification on Trainer]
\begin{center}
\begin{tabular}{ccccc}
\toprule
\textbf{Input $A$} & \textbf{Input $B$} & \textbf{Input $C_{\mathrm{in}}$} & \textbf{Observed Sum LED} & \textbf{Observed Carry LED} \\
\midrule
0 & 0 & 0 & OFF (0) & OFF (0) \\
0 & 0 & 1 & ON (1) & OFF (0) \\
0 & 1 & 0 & ON (1) & OFF (0) \\
0 & 1 & 1 & OFF (0) & ON (1) \\
1 & 0 & 0 & ON (1) & OFF (0) \\
1 & 0 & 1 & OFF (0) & ON (1) \\
1 & 1 & 0 & OFF (0) & ON (1) \\
1 & 1 & 1 & ON (1) & ON (1) \\
\bottomrule
\end{tabular}
\end{center}
\end{observationtable}

\section{Worked Numerical Examples}
\begin{example}[Multi-Bit Binary Addition and Propagation Delay]
A 4-bit ripple-carry adder is constructed by cascading four Full Adders. Each Full Adder has a gate delay of $12\,\mathrm{ns}$ for Sum and $8\,\mathrm{ns}$ for Carry generation. (a) Compute the result of adding $A = 1011_2$ (11) and $B = 0110_2$ (6). (b) Determine the maximum propagation delay required to establish the final valid 4-bit sum and carry-out.
\end{example}
\begin{solutionbox}
(a) Binary addition calculation:
\begin{center}
\begin{tabular}{rccccc}
Carry: & 1 & 1 & 0 & 0 & \\
$A$: & & 1 & 0 & 1 & 1 \\
$B$: & & 0 & 1 & 1 & 0 \\
\midrule
Sum: & 1 & 0 & 0 & 0 & 1 \\
\end{tabular}
\end{center}
The result is $10001_2 = 16 + 1 = 17_{10}$, verifying arithmetic correctness.\\
(b) Propagation Delay:
In a ripple carry adder, the carry must ripple through all four stages:
\[ t_{\mathrm{carry,total}} = 4 \times t_{\mathrm{carry}} = 4 \times 8\,\mathrm{ns} = 32\,\mathrm{ns} \]
The final MSB sum bit requires the final carry plus the XOR sum delay:
\[ t_{\mathrm{total}} = 3 \times t_{\mathrm{carry}} + t_{\mathrm{sum}} = (3 \times 8) + 12 = 24 + 12 = 36\,\mathrm{ns} \]
\end{solutionbox}

\section{Workshop Viva-Voce Questions \& Answers}
\begin{enumerate}[leftmargin=1.5em]
    \item \textbf{Question}: What is the difference between a ripple-carry adder and a carry-lookahead adder?\\
    \textbf{Answer}: In a ripple-carry adder, each Full Adder stage must wait for the carry bit to propagate sequentially from preceding stages, resulting in a total delay proportional to $n$ ($O(n)$). A carry-lookahead adder computes carry bits in parallel across all stages simultaneously using Generate ($G_i = A_i B_i$) and Propagate ($P_i = A_i \oplus B_i$) logic, reducing the delay to a constant $O(1)$ independent of bit width.
    \item \textbf{Question}: How can a binary adder be modified to perform binary subtraction without extra subtractor ICs?\\
    \textbf{Answer}: By utilizing 2's complement arithmetic ($A - B = A + \overline{B} + 1$). Inputs $B$ are passed through XOR gates controlled by a Mode line ($M$). When $M=1$, the XOR gates invert $B$ to produce the 1's complement ($\overline{B}$), and setting the initial carry $C_0 = M = 1$ adds 1, completing the 2's complement subtraction.
\end{enumerate}
